% ======================================================================
\section{作品概述}
% ======================================================================

\subsection{选题背景与意义}

旋转机械、泵阀、风机和传动机构在出现松动、不对中、轴承磨损或叶片损伤时，振动通常比温度、电流等量更早发生变化。故障早期的振动幅度往往只有零点几毫米甚至更小，人在现场很难看出来；若不及时处理，松动会加剧冲击，磨损会扩大间隙，最终可能导致停机、部件飞出或结构疲劳断裂。因此，持续观察振动的频率和幅值，是设备状态监测中最常用、也最有效的手段之一\cite{iso20816}。

目前工程上主要有三类测振方式，各有明显的局限（表~\ref{tab:methods}）。压电或 MEMS 加速度计需要用螺栓、磁座或胶粘固定在被测部位，每个测点占用一路采集通道并要单独布线；对薄板、叶片这类轻质结构，传感器自身的附加质量还会改变被测件的动态特性。激光多普勒测振仪可以非接触测量，精度很高，但一束激光只对准一个点，扫描式设备要逐点移动光斑，多点扫描还需要专用光学和运动控制组件。基于高速相机的数字图像相关（DIC）和运动放大方法可以一次覆盖整片区域\cite{sutton2009,chen2015}，但多依赖高速相机和离线工作站（甚至可以从视频中恢复声音\cite{davis2014}），难以直接放到现场长期运行。

\begin{table}[htbp]
\centering
\caption{常见测振方式与本作品的对比}
\label{tab:methods}
\small
\begin{tabularx}{\linewidth}{@{}L{2.4cm}L{2.3cm}L{2.5cm}YL{2.4cm}@{}}
\toprule
\tabhead \tabh{方式} & \tabh{是否接触} & \tabh{同时测点数} & \tabh{主要限制} & \tabh{部署成本} \\
\midrule
加速度计 & 接触，需固定 & 每通道 1 点 & 逐点安装布线；附加质量影响轻质结构 & 单点低，多点随通道数增长 \\
激光多普勒测振仪 & 非接触 & 1 点（扫描式逐点） & 价格高；对准与反光条件要求严格 & 很高 \\
高速相机 DIC & 非接触 & 整个视场 & 依赖高速相机与离线计算，现场难以常驻 & 高 \\
\textbf{见微（本作品）} & \textbf{非接触} & \textbf{最多 12 点 + 4 个参考点} & 受相机帧率限制，频率上限约为帧率的一半 & \textbf{低（嵌入式 SoC）} \\
\bottomrule
\end{tabularx}
\end{table}

我们的思路是把“相机看振动”这件事做成一台能放在现场的嵌入式仪器：在被测表面贴一块散斑纹理，相机从几十厘米外拍摄，设备在本地完成多点亚像素跟踪、频谱分析和与停机状态的比较，超过阈值就在现场亮灯、鸣响、语音提示并打印报告，同时把结果推送到手机和电脑。整套方案不需要在机器上布线，一台设备可以同时看多个测点，成本以开发板和自制电路为主。

视频前端的数据量决定了这件事必须软硬件协同完成。AN5642 双目模块的每路 OV5640 输出 1280×720 图像，帧率配置为 45 帧/s，按 BGR888 存储时每帧约 2.76 MB，单路约 124 MB/s，两路合计约 249 MB/s；与此同时还要以 1024×600 刷新 HDMI 触摸屏。若完全由 ARM 处理器搬运像素，CPU 时间几乎会被拷贝耗尽。Zynq-7000 恰好提供了 PL 与 PS 共享 DDR 的结构\cite{ug585}：像素级的时序和搬运交给 PL 与 VDMA\cite{pg020}，测量算法、界面和数据管理留在 PS 上用 C++ 实现，便于反复修改和调试。

\subsection{功能与特性}

见微的使用流程很简单：在屏幕或 APP 上用一个矩形框选定测点区域（默认 4 列 3 行共 12 个测点），再在不随设备振动的背景上放置 4 个参考点；电机停机时采集约 15 s 停机基线；启动电机后点击“开始测量”，约 10 s 后得到结果。设备的主要功能如下：

\begin{itemize}
\item \textbf{多点非接触测振：}最多 12 个测点同时做亚像素位移跟踪，4 个背景参考点扣除相机自身的抖动；输出主频、代表测点的峰峰值和有效值，以及逐点的信号/基线比。
\item \textbf{停机基线防误报：}每个测点、每个频点保存停机时的频谱包络，运行时只有超过基线 3 倍的分量才算振动信号，减少环境光闪烁和相机晃动带来的误报。
\item \textbf{声学独立核对：}PL 接收 MP34DT01 PDM 麦克风的原始位流，PS 抽取为 31.25 kHz 音频并做包络谱分析，仅凭声音给出另一条主频估计，与视觉结果交叉核对。
\item \textbf{现场报警与记录：}三色报警灯、蜂鸣器和中文语音分级提示；热敏打印机输出带波形的测量报告；每次测量保存完整 JSON 记录（含 12 个测点、标定、采集质量和声学分析）。
\item \textbf{多端交互：}7 英寸触摸屏提供本地交互，前面板集成旋钮、按键和状态灯；Android、iOS 与 Windows 共用一套 APP，可在局域网直连设备或经 ESP32 热点访问；NFC 卡片可绑定测点编号。
\item \textbf{辅助感知：}PS 端低分辨率双目匹配给出场景深度参考；十字激光辅助对准；预留 IMU、参考加速度计、ToF 测距、温湿度和转速接口。
\end{itemize}

\begin{figure}[htbp]
\centering
\begin{subfigure}[t]{0.54826\linewidth}\centering\vspace{0pt}
\includegraphics[height=50.02mm]{\figpath/live/full_alarm.jpg}
\caption{2026-11-04 板端实测现场：左侧为减速电机带动的散斑铝片，右侧为见微整机，报警灯已点亮}
\end{subfigure}\hfill
\begin{subfigure}[t]{0.41174\linewidth}\centering\vspace{0pt}
\includegraphics[height=50.02mm]{\figpath/photo/device_lit.jpg}
\caption{整机前视：双目相机居中，两侧为补光灯，下方十字激光用于对准}
\end{subfigure}
\caption{见微整机与实测场景}
\label{fig:scene}
\end{figure}

\subsection{应用领域}

见微面向“能看见、但不方便贴传感器”的振动观测场景：实验室和教学中的电机、薄板、悬臂梁振动演示；产线上小型电机、风机、泵的巡检式趋势监测；轻质结构（风扇叶片、薄壁罩壳、PCB 板件）的非接触测量；以及需要同时比较多个位置振动差异的场合，例如判断振动从哪个部位传来。由于只需在被测面贴散斑纹理、在背景中保留静止参考，见微也适合临时布点的现场排查。

当前版本的频率上限受 45 帧/s 的采样率限制，约为 22.5 Hz，适合 1350 r/min 以下的一倍转频观测；更高频率的振动需要更高帧率的相机。

\subsection{主要技术特点}

\paragraph{按带宽与时序划分软硬件。}PL 只做“必须准时”的事：两路相机的像素接收与帧缓存、HDMI 时序、2 MHz PDM 时钟与采样、20 kHz 补光 PWM 与心跳看门狗、打印串口的 CTS 流控。PS 做“需要灵活修改”的事：跟踪、频谱、判定、界面和记录。ESP32 只承担辅助外设，不参与测振结论。

\paragraph{面向现场的鲁棒判定。}跟踪采用逆组合 Lucas–Kanade 方法，模板梯度和 Hessian 只算一次；每帧给出相关系数、结构张量条件数等质量指标。背景参考点用中位数加 0.25 像素内点平均估计相机平移；分析窗用帧序号和时间戳检查丢帧和抖动，只在缺帧比例不超过 5\% 时重采样；判定以逐点停机基线包络为门限。

\paragraph{声、像两条独立链路。}视觉测得的是位移，声学测得的是机器噪声随转动周期的调制，两者物理来源不同。两条链路给出一致的主频，比任何一条单独给出的结果更可信；不一致时，系统会标注谐波关系或差异，提示检查。

\paragraph{结果可追溯。}每次测量记录相机实测帧率、丢帧、跟踪质量、标定编号、阈值快照和声学分析，支持结果回看、报告打印和多端查看。

\subsection{主要性能指标}

表~\ref{tab:perf} 的实测数值来自 2026-11-04 的板端测量记录 \code{jianwei-ax7010-20261104-020520}，资源与时序来自 Vivado 2025.2 实现报告。

\begin{table}[htbp]
\centering
\caption{主要性能指标}
\label{tab:perf}
\small
\begin{tabularx}{\linewidth}{@{}L{3.6cm}L{4.6cm}Y@{}}
\toprule
\tabhead \tabh{项目} & \tabh{指标} & \tabh{说明} \\
\midrule
图像输入 & 双路 OV5640，1280×720 & AN5642 模块，测振使用相机 1 \\
采样率 & 45.0477 帧/s，抖动 0.417 ms & 由帧序号与单调时钟时间戳实测 \\
测点数量 & 最多 12 个测点 + 4 个参考点 & 默认 4 列 3 行，模板 25×25 像素 \\
频率范围 & 0.5 Hz～22.5 Hz & 45 帧/s 视觉采样，用于低频振动与转频观测 \\
频率分辨率 & 窗口频点间隔 0.352 Hz & 128 帧窗，峰值插值后取 11 窗中位数 \\
理论转频对照 & 1200 r/min，对应 20 Hz & 视觉实测 19.9476 Hz，相对差约 0.26\% \\
主频一致性 & 视觉 19.9476 Hz，声学 19.9277 Hz & 差 0.0199 Hz，约 0.10\% \\
位移结果 & 峰峰值 0.2090 mm，有效值 0.0701 mm & 试验尺标 0.20309 mm/像素 \\
位移噪声 & 约 0.005～0.011 像素（约 1.1～2.3 µm） & 主频处等效背景谱幅，由基线比推算，见 \ref{sec:noise} 节 \\
有效性 & 12/12 测点，11/11 窗通过 & 逐点信号/基线比 36～95，门限 3 \\
PL 资源 & LUT 50.8\%，FF 41.0\%，BRAM 50.8\% & XC7Z010，Vivado 2025.2 实现后 \\
时序 & WNS +0.756 ns，WHS +0.028 ns & 已约束路径全部满足 \\
声学通道 & 2 MHz PDM → 31.25 kHz PCM & 滤波器计算响应：通带波动 ±0.004 dB \\
整机功耗 & 约 9.3 W（12 V 输入） & 板载 INA226 读数，测量时补光 60\% \\
\bottomrule
\end{tabularx}
\end{table}
\FloatBarrier


\subsection{整机实测结果}

见微已完成从图像采集、测振分析到声光报警、语音、打印和 APP 查看的一体化运行。本次电机转速为 1200 r/min，理论转频为 $f=n/60=20$ Hz。视觉实测主频为 19.9476 Hz，与理论转频相差约 0.0524 Hz，相对差约 0.26\%，两者十分接近。峰峰值为 0.2090 mm，452 帧采集无丢帧，报警与测量报告在现场同步输出。

2026 年 11 月 4 日 02:05 的完整测量如表~\ref{tab:record} 所示。被测物为减速电机驱动的散斑铝片，相机 1 输出 1280×720 图像，测量框位于 (197, 235)，大小 256×124 像素，12 个测点按 4 列 3 行排布，4 个参考点贴在铝片下方静止的背景上（图~\ref{fig:points}）。先采 15 s 停机基线，再启动电机测量约 10 s。

\begin{table}[htbp]
\centering
\caption{2026-11-04 板端测量记录 \code{jianwei-ax7010-20261104-020520}}
\label{tab:record}
\small
\begin{tabularx}{\linewidth}{@{}L{3.6cm}L{4.3cm}Y@{}}
\toprule
\tabhead \tabh{项目} & \tabh{结果} & \tabh{说明} \\
\midrule
电机转速 / 理论转频 & 1200 r/min / 20 Hz & $f=n/60$，实测主频相对差约 0.26\% \\
开始时间 / 时长 & 02:05:10.6 / 10.012 s & RTC 时间源，单调时钟计时 \\
帧数 / 帧率 & 452 帧 / 45.0477 帧/s & 收到帧率与相机帧率相同 \\
丢帧 / 跟踪失败 / 饱和 & 0 / 0 / 0 & 无重采样补点 \\
帧间隔抖动 & 0.417 ms & 约为帧周期的 1.9\% \\
分析窗 & 11 个，全部为信号窗 & 128 帧窗，每 32 帧一次 \\
视觉主频 & 19.9476 Hz & 11 窗范围 19.851～19.964 Hz \\
声学主频 & 19.9277 Hz & 3000～8000 Hz 包络，27.2 dB \\
主频差 & 0.0199 Hz（0.10\%） & 容差 0.399 Hz，判为一致 \\
超过基线的测点 & 12/12 & 信号/基线比 36.3～95.0 \\
代表测点 & 第 2 点（比值 95.0） & 波形、频谱检查通过 \\
峰峰值 & 1.0294 像素 / 0.2090 mm & 代表测点 $y$ 方向时域值 \\
有效值 & 0.3451 像素 / 0.0701 mm & 同一序列，去线性趋势 \\
尺标比例 & 0.20309 mm/像素 & 试验尺标，$y$ 方向 \\
判级 & 报警 & 峰峰值 0.209 mm，高于报警阈值 0.02 mm \\
照明 & 开启，60\% & 十字激光关闭 \\
\bottomrule
\end{tabularx}
\end{table}
\FloatBarrier

帧率按首末帧之间的 451 个间隔计算：$f_s=(452-1)/10.01160=45.0477$ 帧/s。记录保留计算精度，屏幕、打印和视频以 19.95 Hz、19.93 Hz、0.209 mm 显示视觉主频、声学主频和峰峰值。


\begin{figure}[htbp]
\centering
\begin{subfigure}[b]{0.47\linewidth}\centering
\begin{tikzpicture}
\begin{axis}[width=\linewidth,height=6.4cm,xmin=180,xmax=470,ymin=220,ymax=545,y dir=reverse,
  axis equal image,xlabel={$x$ / 像素},ylabel={$y$ / 像素},tick label style={font=\scriptsize},label style={font=\footnotesize},
  colormap={warm}{rgb255(0)=(255,236,179) rgb255(1)=(242,140,40) rgb255(2)=(178,34,34)},
  colorbar,colorbar style={width=2.5mm,ylabel={信号/基线比},ylabel style={font=\scriptsize},tick label style={font=\scriptsize}},
  point meta min=30,point meta max=100]
\draw[navy,thick] (axis cs:197,235) rectangle (axis cs:453,359);
\addplot[scatter,only marks,mark=*,mark size=3.6pt,scatter src=explicit] coordinates {
 (229,255.67)[81.1] (293,255.67)[95.0] (357,255.67)[67.3] (421,255.67)[59.7]
 (229,297)[75.6] (293,297)[62.3] (357,297)[60.7] (421,297)[36.3]
 (229,338.33)[65.0] (293,338.33)[67.9] (357,338.33)[54.1] (421,338.33)[37.3]};
\addplot[only marks,mark=square*,mark size=3pt,color=navy] coordinates {(266,494) (298,494) (266,526) (298,526)};
\node[font=\scriptsize,anchor=south west,text=navy] at (axis cs:197,232) {测量框};
\node[font=\scriptsize,anchor=west,text=navy] at (axis cs:305,510) {参考点};
\end{axis}
\end{tikzpicture}
\caption{测点与参考点在图像中的位置}
\end{subfigure}\hfill
\begin{subfigure}[b]{0.5\linewidth}\centering
\begin{tikzpicture}
\begin{axis}[width=\linewidth,height=6.4cm,ybar,bar width=8pt,ymin=0,ymax=105,xmin=0.3,xmax=12.7,
  xtick={1,...,12},xlabel={测点编号},ylabel={信号/基线比 $g_i$},tick label style={font=\scriptsize},
  label style={font=\footnotesize},nodes near coords,nodes near coords style={font=\tiny,rotate=90,anchor=west},
  grid=major,grid style={draw=gray!20}]
\addplot[fill=amber!80,draw=amber!60!black] coordinates {(1,81.1) (2,95.0) (3,67.3) (4,59.7) (5,75.6) (6,62.3) (7,60.7) (8,36.3) (9,65.0) (10,67.9) (11,54.1) (12,37.3)};
\draw[axred,thick,dashed] (axis cs:0.3,3) -- (axis cs:12.7,3);
\node[font=\scriptsize,text=axred,anchor=south east] at (axis cs:12.6,4) {门限 3};
\end{axis}
\end{tikzpicture}
\caption{逐点信号/基线比}
\end{subfigure}
\caption{2026-11-04 实测的测点布局与判定结果}
\label{fig:points}
\end{figure}

12 个测点的信号/基线比都在 36 以上，远高于门限 3。测量框右列的测点 4、8、12 振幅较小（谐波有效值 0.26～0.30 像素），比值最低；左列和中部测点振幅较大，比值最高，铝片各处振幅的差异在测点之间清楚地表现出来。视觉与声学两条独立链路给出的主频相差 0.02 Hz，远小于视觉单窗 0.352 Hz 的频点间隔；11 个窗口的主频分布在 19.851～19.964 Hz 之间，取中位数后结果稳定。

从按下“开始测量”到报警输出的过程与设计流程一致：基线采集完成时播放“停机基线已采集，请开始测量”；启动电机后点击开始，10 s 测量结束即给出结果，峰峰值 0.209 mm 超过 0.02 mm 报警阈值，红灯点亮、蜂鸣器间歇鸣响并播报“振动异常，请及时检查设备”，手机 APP 同时弹出报警；随后从结果页触发打印，热敏打印机输出含主频、峰峰值、12/12 测点和波形的报告（图~\ref{fig:uilive}、图~\ref{fig:app}）。


\begin{figure}[htbp]
\centering
\begin{subfigure}[t]{0.48000\linewidth}\centering\vspace{0pt}
\includegraphics[height=43.73mm]{\figpath/live/scr_meas.jpg}
\caption{测量中}
\end{subfigure}\hfill
\begin{subfigure}[t]{0.48000\linewidth}\centering\vspace{0pt}
\includegraphics[height=43.73mm]{\figpath/live/scr_result.jpg}
\caption{结果页：19.95 Hz，12/12，0.209 mm}
\end{subfigure}
\par\vspace{6mm}
\begin{subfigure}[t]{0.96000\linewidth}\centering\vspace{0pt}
\includegraphics[height=80.00mm]{\figpath/live/scr_receipt.jpg}
\caption{打印出的测量报告}
\end{subfigure}
\caption{2026-11-04 实测时的屏幕与打印输出}
\label{fig:uilive}
\end{figure}
\FloatBarrier
\begin{figure}[htbp]
\centering
\begin{subfigure}[t]{0.48000\linewidth}\centering\vspace{0pt}
\includegraphics[height=90.00mm]{\figpath/live/app_alarm.jpg}
\caption{实测报警推送}
\end{subfigure}\hfill
\begin{subfigure}[t]{0.48000\linewidth}\centering\vspace{0pt}
\includegraphics[height=90.00mm]{\figpath/live/app_result.jpg}
\caption{实测结果详情}
\end{subfigure}
\caption{APP 实测报警推送与结果详情}
\label{fig:app}
\end{figure}
\FloatBarrier


\subsection{主要创新点}

\begin{enumerate}[leftmargin=2em]
\item \textbf{一台相机同时测多个点，不接触、不布线。}在 Zynq-7010 这样资源有限的器件上实现 12 点亚像素跟踪与逐点频谱判定，并把采集、分析、报警、打印和移动端推送集成在一台独立设备中。
\item \textbf{背景参考点与停机基线组合的现场判定方法。}用多个静止参考点的稳健平移估计扣除相机抖动，再以逐点、逐频点的停机频谱包络作为门限，抑制环境光闪烁和设备晃动带来的误报；判定只依赖本次机位的数据，不需要事先知道转速。
\item \textbf{PL 接收 PDM 麦克风，用声音独立核对振动频率。}PL 只保留 2 MHz 时钟、相位采样、FIFO 与时间戳，抽取滤波放到 PS 并与参考模型逐位一致；声学包络谱给出与视觉无关的主频估计，实测两者相差 0.1\%。
\item \textbf{资源驱动的软硬件划分与可追溯的优化过程。}把双目深度从 PL 计算核迁到 PS 辅助路径，整机实现净减少 5081 个 LUT 和 11 块 BRAM，换取视频、声学与外设链路的稳定；资源配置随整机功能逐步调整。
\item \textbf{自研载板、前面板与外壳。}主载板集成 12 V 电源树、功率监测、掉电预警、RTC、IO 扩展、跨电源域缓冲和全部外设接口；外壳将相机支架下倾约 25°，集成补光、激光、打印机与报警灯。
\end{enumerate}

\subsection{设计流程}

项目从 2026 年 9 月下旬选题开始，按“需求与约束—架构划分—电路与 PL—PS 算法—外设联调—整机装配—板端实测”的顺序推进，各阶段之间多次回溯（图~\ref{fig:timeline}）。每一轮修改都先固定版本和输入，再按单元测试、主机集成、ARM 交叉编译、必要的 Vivado 实现、板端运行和现场确认逐层验证；发现问题时先划清是采集、传输、计算还是显示环节，再只改变一个变量复测（图~\ref{fig:flow}）。

\begin{figure}[htbp]
\centering
\begin{tikzpicture}[x=1cm,y=1cm]
\draw[line width=1.6pt,navy] (0.2,0) -- (15.4,0);
\foreach \x/\d/\t/\c [count=\i] in {
  0.6/9 月下旬/{选题与\\需求分析}/navy,
  2.45/10-19/{两块电路板\\原理图定稿}/amber,
  4.3/10-20/{PCB 投板}/amber,
  6.15/10-21/{PL 双目核与\\LK 累加核验证}/axred,
  8.0/10-22/{PL、PS 与 ESP32\\分工定稿}/navy,
  9.85/10-25/{程序冻结\\多点核接入}/navy,
  11.7/10-31/{Nano 0.2.1\\报警标定}/espc!55!black,
  13.55/11-02/{Nano 0.2.3\\13 条语音}/espc!55!black,
  15.0/11-04/{整机\\板端实测}/axred}{
  \fill[\c] (\x,0) circle (3.2pt);
  \ifodd\i
    \node[above=3pt,font=\footnotesize\bfseries,text=\c] at (\x,0) {\d};
    \node[below=4pt,align=center,font=\footnotesize\linespread{1.05}\selectfont] at (\x,0) {\t};
  \else
    \node[below=3pt,font=\footnotesize\bfseries,text=\c] at (\x,0) {\d};
    \node[above=4pt,align=center,font=\footnotesize\linespread{1.05}\selectfont] at (\x,0) {\t};
  \fi
}
\end{tikzpicture}
\caption{开发时间线}
\label{fig:timeline}
\end{figure}

\begin{figure}[htbp]
\centering
\begin{tikzpicture}[x=1cm,y=1cm]
\node[blk=psc,minimum width=2.6cm] (req) at (0,0) {需求与约束\\验收条件};
\node[blk=psc,minimum width=2.6cm] (ev) at (3.5,0) {证据采集\\源码、日志、读数};
\node[blk=psc,minimum width=2.6cm] (hy) at (7,0) {可证伪的假设\\划清环节};
\node[blk=psc,minimum width=2.6cm] (fix) at (10.5,0) {最小修改\\只改一个变量};
\node[blk=ownc,minimum width=13.1cm] (ver) at (5.25,-1.7) {分层验证：单元测试 → 主机集成 → 交叉编译 → Vivado 实现 → 板端运行 → 现场确认};
\node[blk=espc,minimum width=2.6cm] (rec) at (1.75,-3.4) {结果回写\\版本、证据、未决项};
\node[blk=espc,minimum width=2.6cm] (sk) at (8.75,-3.4) {经验提炼\\技能包};
\draw[arr] (req) -- (ev);
\draw[arr] (ev) -- (hy);
\draw[arr] (hy) -- (fix);
\draw[arr] (fix) -- (fix |- ver.north);
\draw[arr,dashed] (ev |- ver.north) -- node[right,lbl]{未通过} (ev);
\draw[arr] (rec |- ver.south) -- (rec);
\draw[arr] (rec) -- node[above,lbl]{重复出现的方法} (sk);
\end{tikzpicture}
\caption{每一轮迭代的证据闭环}
\label{fig:flow}
\end{figure}

\FloatBarrier
